% GENERATED FILE. Edit canonical lessons, then run npm run generate:books.
% canonical-source-hash: fnv1a64:40c72d7025304e25
% canonical-lessons: GE-C485-nacken, GE-C485-ellbogen, GE-C485-handgelenk, GE-C485-knoechel, GE-C485-huefte

\chapter{Nacken, Ellbogen und Knöchel}
\label{ch:ge-nacken-ellbogen-und-knochel}
\hlchaptermodality{485}

\begin{chapteropening}
\textbf{By the end of this chapter:} \emph{I can name the joints and describe aches.}

Complete the last lesson of chapter 485: I can name the joints and describe aches.
\end{chapteropening}

\section[der Nacken]{der Nacken — back of the neck, nape}

\label{lesson:GE-C485-nacken}



\begin{quote}
Before the new one: say the German for eyebrow, then the German for eyelash.
\end{quote}

\subsection*{You'll want to know: der Nacken}
\textbf{der Nacken} — \textquotedblleft{}back of the neck, nape\textquotedblright{}.

\begin{quote}
Nach der Arbeit am Computer tut mir der Nacken weh.
\end{quote}

\begin{quote}
\emph{After working at the computer my neck hurts.}
\end{quote}

Plural: \emph{die Nacken}. \emph{Der Hals} is the whole neck; \emph{der Nacken} is the back of it.

\begin{cousinweb}
\emph{Nacken} is inherited from Middle High German \emph{nacke}, from Proto-Germanic \emph{hnakkô} \textquotedblleft{}nape, neck\textquotedblright{} (EtymDB).
\end{cousinweb}

\begin{grammarlens}[title={What you've built}]
One more word for telling, explaining and describing at B1.
\end{grammarlens}

\subsection*{Your turn}
\begin{itemize}
\raggedright
  \item \emph{Say it:} \emph{der Nacken}
  \item \emph{Say it:} \emph{der Nacken}, once more
  \item \emph{Say it:} a sentence of your own with \emph{der Nacken}
  \item \emph{Recall:} say the German for eyebrow, then the German for eyelash, then say \emph{der Nacken} again
\end{itemize}

\begin{tcolorbox}[breakable,colback=teal!4,colframe=teal!35!black,title={Before you move on}]
What does der Nacken mean? (\textquotedblleft{}Back of the neck, nape\textquotedblright{}.) Say it once more.
\end{tcolorbox}
\section[der Ellbogen]{der Ellbogen — elbow}

\label{lesson:GE-C485-ellbogen}



\begin{quote}
Before the new one: say the German for eyelash, then the German for back of the neck.
\end{quote}

\subsection*{You'll want to know: der Ellbogen}
\textbf{der Ellbogen} — \textquotedblleft{}elbow\textquotedblright{}.

\begin{quote}
Er hat sich beim Sturz den Ellbogen verletzt.
\end{quote}

\begin{quote}
\emph{He hurt his elbow in the fall.}
\end{quote}

Plural: \emph{die Ellbogen}. Also spelt \emph{der Ellenbogen}.

\begin{cousinweb}
\emph{Ellbogen} is inherited from Middle High German \emph{elenboge} and Old High German \emph{elinbogo}, from Proto-Germanic \emph{alinabugô} \textquotedblleft{}elbow\textquotedblright{} (EtymDB): the bend of the forearm.
\end{cousinweb}

\begin{grammarlens}[title={What you've built}]
One more word for telling, explaining and describing at B1.
\end{grammarlens}

\subsection*{Your turn}
\begin{itemize}
\raggedright
  \item \emph{Say it:} \emph{der Ellbogen}
  \item \emph{Say it:} \emph{der Ellbogen}, once more
  \item \emph{Say it:} a sentence of your own with \emph{der Ellbogen}
  \item \emph{Recall:} say the German for eyelash, then the German for back of the neck, then say \emph{der Ellbogen} again
\end{itemize}

\begin{tcolorbox}[breakable,colback=teal!4,colframe=teal!35!black,title={Before you move on}]
What does der Ellbogen mean? (\textquotedblleft{}Elbow\textquotedblright{}.) Say it once more.
\end{tcolorbox}
\section[das Handgelenk]{das Handgelenk — wrist}

\label{lesson:GE-C485-handgelenk}



\begin{quote}
Before the new one: say the German for back of the neck, then the German for elbow.
\end{quote}

\subsection*{You'll want to know: das Handgelenk}
\textbf{das Handgelenk} — \textquotedblleft{}wrist\textquotedblright{}.

\begin{quote}
Sie trägt eine Uhr am linken Handgelenk.
\end{quote}

\begin{quote}
\emph{She wears a watch on her left wrist.}
\end{quote}

Plural: \emph{die Handgelenke}. \emph{Das Gelenk} — joint.

\begin{cousinweb}
\emph{Handgelenk} is built on \emph{Hand} and \emph{Gelenk} \textquotedblleft{}joint\textquotedblright{}, inherited from Middle High German \emph{gelenke} and Old High German \emph{gilenki} \textquotedblleft{}joint, bend, waist\textquotedblright{} (EtymDB).
\end{cousinweb}

\begin{grammarlens}[title={What you've built}]
One more word for telling, explaining and describing at B1.
\end{grammarlens}

\subsection*{Your turn}
\begin{itemize}
\raggedright
  \item \emph{Say it:} \emph{das Handgelenk}
  \item \emph{Say it:} \emph{das Handgelenk}, once more
  \item \emph{Say it:} a sentence of your own with \emph{das Handgelenk}
  \item \emph{Recall:} say the German for back of the neck, then the German for elbow, then say \emph{das Handgelenk} again
\end{itemize}

\begin{tcolorbox}[breakable,colback=teal!4,colframe=teal!35!black,title={Before you move on}]
What does das Handgelenk mean? (\textquotedblleft{}Wrist\textquotedblright{}.) Say it once more.
\end{tcolorbox}
\section[der Knöchel]{der Knöchel — ankle; knuckle}

\label{lesson:GE-C485-knoechel}



\begin{quote}
Before the new one: say the German for elbow, then the German for wrist.
\end{quote}

\subsection*{You'll want to know: der Knöchel}
\textbf{der Knöchel} — \textquotedblleft{}ankle; knuckle\textquotedblright{}.

\begin{quote}
Beim Wandern habe ich mir den Knöchel verstaucht.
\end{quote}

\begin{quote}
\emph{I sprained my ankle while hiking.}
\end{quote}

Plural: \emph{die Knöchel}. Also a knuckle of the finger.

\begin{cousinweb}
\emph{Knöchel} is usually traced to a small form of \emph{Knochen} \textquotedblleft{}bone\textquotedblright{}, from Middle High German \emph{knöchel}; EtymDB lists the German word without a source.
\end{cousinweb}

\begin{grammarlens}[title={What you've built}]
One more word for telling, explaining and describing at B1.
\end{grammarlens}

\subsection*{Your turn}
\begin{itemize}
\raggedright
  \item \emph{Say it:} \emph{der Knöchel}
  \item \emph{Say it:} \emph{der Knöchel}, once more
  \item \emph{Say it:} a sentence of your own with \emph{der Knöchel}
  \item \emph{Recall:} say the German for elbow, then the German for wrist, then say \emph{der Knöchel} again
\end{itemize}

\begin{tcolorbox}[breakable,colback=teal!4,colframe=teal!35!black,title={Before you move on}]
What does der Knöchel mean? (\textquotedblleft{}Ankle; knuckle\textquotedblright{}.) Say it once more.
\end{tcolorbox}
\section[die Hüfte]{die Hüfte — hip}

\label{lesson:GE-C485-huefte}



\begin{quote}
Before the new one: say the German for wrist, then the German for ankle.
\end{quote}

\subsection*{You'll want to know: die Hüfte}
\textbf{die Hüfte} — \textquotedblleft{}hip\textquotedblright{}.

\begin{quote}
Meine Oma bekommt eine neue Hüfte.
\end{quote}

\begin{quote}
\emph{My grandma is getting a new hip.}
\end{quote}

Plural: \emph{die Hüften}.

\begin{cousinweb}
\emph{Hüfte} is usually traced to Old High German \emph{huf}, which EtymDB derives from Proto-Germanic \emph{hupiz} \textquotedblleft{}hip\textquotedblright{}; EtymDB sets \emph{Hüfte} beside English \emph{hip}, a cousin.
\end{cousinweb}

\begin{grammarlens}[title={What you've built}]
One more word for telling, explaining and describing at B1.
\end{grammarlens}

\subsection*{Your turn}
\begin{itemize}
\raggedright
  \item \emph{Say it:} \emph{die Hüfte}
  \item \emph{Say it:} \emph{die Hüfte}, once more
  \item \emph{Say it:} a sentence of your own with \emph{die Hüfte}
  \item \emph{Recall:} say the German for wrist, then the German for ankle, then say \emph{die Hüfte} again
\end{itemize}

\begin{tcolorbox}[breakable,colback=teal!4,colframe=teal!35!black,title={Before you move on}]
What does die Hüfte mean? (\textquotedblleft{}Hip\textquotedblright{}.) Say it once more.
\end{tcolorbox}
